\section{Threats and Limitations}
\label{sec:threats}
In this work, we focus on automated repair of bugs drawn from Google's internal repository.
Other industrial settings may have bugs that reflect different characteristics. However,
Google's codebase is likely to share properties with other large software providers.
To create our GITS-Eval dataset we relied on manual assessment of properties that could
not be automatically enforced (see Section~\ref{sec:challengeset}). To mitigate this risk,
we employed multiple annotators.


Importantly, we present Passerine as a proof-of-concept for agent-based APR in an industrial setting.
As a result, we make no claims of how Passerine would fare on other benchmarks such as SWE-Bench. Passerine's
design is also purposefully simple and it is possible improvements in both plausible and valid patch generation
rates would come from a more-sophisticated design. Assessing the viability of such improvements remains part of 
our future work. Finally, we note that patch validity was assessed manually and is subject to the usual
challenges of such evaluation. To mitigate this risk we employed multiple annotators for cases that
were considered ambiguous.

